% GENERATED FILE. Edit canonical lessons, then run npm run generate:books.
% canonical-source-hash: fnv1a64:ed6d80c2dbefbce4
% canonical-lessons: MW-C257-ris, MW-C257-rusno, MW-C257-laj, MW-C257-nisaso, MW-C257-olmo

\chapter{Anger and Shame}
\label{ch:mw-anger-and-shame}
\hlchaptermodality{257}

\begin{chapteropening}
\textbf{By the end of this chapter:} \emph{I can say when someone is angry or shy.}

Complete the last lesson of chapter 257: I can say when someone is angry or shy.
\end{chapteropening}

\section[\texorpdfstring{\mw{रीस} (rīs)}{रीस (rīs)}]{\mw{रीस} (rīs) — anger}

\label{lesson:MW-C257-ris}



\begin{quote}
Before the new one: say the Marwadi for joy, then the Marwadi for high spirits.
\end{quote}

\subsection*{You'll want to know: \mw{रीस}}
\textbf{\mw{रीस}} — \emph{rīs} — \textquotedblleft{}anger\textquotedblright{}.

\begin{quote}
\mw{रीस} \mw{मत} \mw{करो।}
\end{quote}

\begin{quote}
\emph{rīs mat karo.}
\end{quote}

\begin{quote}
Don't get angry.
\end{quote}

A feminine noun.

\begin{cousinweb}
\textbf{\mw{रीस}} (\emph{rīs}) is inherited from Sanskrit \emph{riṣ}, \textquotedblleft{}injury\textquotedblright{}. CDIAL gives no Marwari form in this entry; it lists Gujarati \emph{rīs} and Hindi \emph{ris}, both \textquotedblleft{}anger\textquotedblright{}, under it; how harm became anger is not spelled out there (CDIAL 10746).
\end{cousinweb}

\begin{grammarlens}[title={What you've built}]
One more word for this chapter.
\end{grammarlens}

\subsection*{Your turn}
\begin{itemize}
\raggedright
  \item \emph{Say it:} \emph{rīs}
  \item \emph{Say it:} \emph{rīs}, once more
  \item \emph{Say it:} \emph{rīs mat karo.} — Don't get angry.
  \item \emph{Recall:} say the Marwadi for joy, then the Marwadi for high spirits, then say \emph{rīs} again
\end{itemize}

\begin{tcolorbox}[breakable,colback=teal!4,colframe=teal!35!black,title={Before you move on}]
What does \mw{रीस} mean? (\textquotedblleft{}Anger\textquotedblright{}.) Say it once more.
\end{tcolorbox}
\section[\texorpdfstring{\mw{रूसणो} (rūsṇo)}{रूसणो (rūsṇo)}]{\mw{रूसणो} (rūsṇo) — to sulk; to take offence}

\label{lesson:MW-C257-rusno}



\begin{quote}
Before the new one: say the Marwadi for high spirits, then the Marwadi for anger.
\end{quote}

\subsection*{You'll want to know: \mw{रूसणो}}
\textbf{\mw{रूसणो}} — \emph{rūsṇo} — \textquotedblleft{}to sulk; to take offence\textquotedblright{}.

\begin{quote}
\mw{छोटी} \mw{बहन} \mw{रोज} \mw{रूसै} \mw{है।}
\end{quote}

\begin{quote}
\emph{chhoṭī bahan roj rūsai hai.}
\end{quote}

\begin{quote}
My little sister sulks every day.
\end{quote}

Past (he …): \mw{वो} \mw{बेटो} \mw{रूस्यो}, \emph{vo beṭo rūsyo}.

\begin{cousinweb}
\textbf{\mw{रूसणो}} (\emph{rūsṇo}) is inherited from Sanskrit \emph{ruṣyati}, \textquotedblleft{}is angry\textquotedblright{}; CDIAL lists Old Marwari \emph{rūsaï}, \textquotedblleft{}to get angry\textquotedblright{}, under it (CDIAL 10794).
\end{cousinweb}

\begin{grammarlens}[title={What you've built}]
One more word for this chapter.
\end{grammarlens}

\subsection*{Your turn}
\begin{itemize}
\raggedright
  \item \emph{Say it:} \emph{rūsṇo}
  \item \emph{Say it:} \emph{rūsṇo}, once more
  \item \emph{Say it:} \emph{chhoṭī bahan roj rūsai hai.} — My little sister sulks every day.
  \item \emph{Recall:} say the Marwadi for high spirits, then the Marwadi for anger, then say \emph{rūsṇo} again
\end{itemize}

\begin{tcolorbox}[breakable,colback=teal!4,colframe=teal!35!black,title={Before you move on}]
What does \mw{रूसणो} mean? (\textquotedblleft{}To sulk; to take offence\textquotedblright{}.) Say it once more.
\end{tcolorbox}
\section[\texorpdfstring{\mw{लाज} (lāj)}{लाज (lāj)}]{\mw{लाज} (lāj) — modesty; shame}

\label{lesson:MW-C257-laj}



\begin{quote}
Before the new one: say the Marwadi for anger, then the Marwadi for to sulk.
\end{quote}

\subsection*{You'll want to know: \mw{लाज}}
\textbf{\mw{लाज}} — \emph{lāj} — \textquotedblleft{}modesty; shame\textquotedblright{}.

\begin{quote}
\mw{छोरी} \mw{घणी} \mw{लाज} \mw{करै} \mw{है।}
\end{quote}

\begin{quote}
\emph{chhorī ghaṇī lāj karai hai.}
\end{quote}

\begin{quote}
The girl is very modest.
\end{quote}

A feminine noun.

\begin{cousinweb}
\textbf{\mw{लाज}} (\emph{lāj}) is inherited from Sanskrit \emph{lajjā}, \textquotedblleft{}shame, modesty\textquotedblright{}; CDIAL lists Marwari \emph{lāj} under it (CDIAL 10910).
\end{cousinweb}

\begin{grammarlens}[title={What you've built}]
One more word for this chapter.
\end{grammarlens}

\subsection*{Your turn}
\begin{itemize}
\raggedright
  \item \emph{Say it:} \emph{lāj}
  \item \emph{Say it:} \emph{lāj}, once more
  \item \emph{Say it:} \emph{chhorī ghaṇī lāj karai hai.} — The girl is very modest.
  \item \emph{Recall:} say the Marwadi for anger, then the Marwadi for to sulk, then say \emph{lāj} again
\end{itemize}

\begin{tcolorbox}[breakable,colback=teal!4,colframe=teal!35!black,title={Before you move on}]
What does \mw{लाज} mean? (\textquotedblleft{}Modesty; shame\textquotedblright{}.) Say it once more.
\end{tcolorbox}
\section[\texorpdfstring{\mw{निसासो} (nisāso)}{निसासो (nisāso)}]{\mw{निसासो} (nisāso) — a sigh}

\label{lesson:MW-C257-nisaso}



\begin{quote}
Before the new one: say the Marwadi for to sulk, then the Marwadi for modesty.
\end{quote}

\subsection*{You'll want to know: \mw{निसासो}}
\textbf{\mw{निसासो}} — \emph{nisāso} — \textquotedblleft{}a sigh\textquotedblright{}.

\begin{quote}
\mw{दादी} \mw{निसासो} \mw{नाखै} \mw{है।}
\end{quote}

\begin{quote}
\emph{dādī nisāso nākhai hai.}
\end{quote}

\begin{quote}
Grandmother heaves a sigh.
\end{quote}

A masculine noun.

\begin{cousinweb}
\textbf{\mw{निसासो}} (\emph{nisāso}) is inherited from Sanskrit \emph{niḥśvāsa}, \textquotedblleft{}a sigh\textquotedblright{}; CDIAL lists Marwari \emph{nisāso} under it (CDIAL 7111).
\end{cousinweb}

\begin{grammarlens}[title={What you've built}]
One more word for this chapter.
\end{grammarlens}

\subsection*{Your turn}
\begin{itemize}
\raggedright
  \item \emph{Say it:} \emph{nisāso}
  \item \emph{Say it:} \emph{nisāso}, once more
  \item \emph{Say it:} \emph{dādī nisāso nākhai hai.} — Grandmother heaves a sigh.
  \item \emph{Recall:} say the Marwadi for to sulk, then the Marwadi for modesty, then say \emph{nisāso} again
\end{itemize}

\begin{tcolorbox}[breakable,colback=teal!4,colframe=teal!35!black,title={Before you move on}]
What does \mw{निसासो} mean? (\textquotedblleft{}A sigh\textquotedblright{}.) Say it once more.
\end{tcolorbox}
\section[\texorpdfstring{\mw{ओळमो} (oḷmo)}{ओळमो (oḷmo)}]{\mw{ओळमो} (oḷmo) — a reproach; a complaint}

\label{lesson:MW-C257-olmo}



\begin{quote}
Before the new one: say the Marwadi for modesty, then the Marwadi for a sigh.
\end{quote}

\subsection*{You'll want to know: \mw{ओळमो}}
\textbf{\mw{ओळमो}} — \emph{oḷmo} — \textquotedblleft{}a reproach; a complaint\textquotedblright{}.

\begin{quote}
\mw{यो} \mw{ओळमो} \mw{साचो} \mw{है।}
\end{quote}

\begin{quote}
\emph{yo oḷmo sācho hai.}
\end{quote}

\begin{quote}
This complaint is fair.
\end{quote}

A masculine noun.

\begin{cousinweb}
\textbf{\mw{ओळमो}} (\emph{oḷmo}) is inherited from Sanskrit \emph{upālambha}, \textquotedblleft{}reproach\textquotedblright{}; CDIAL lists Marwari \emph{olamũ}, \textquotedblleft{}reproach\textquotedblright{}, under it (CDIAL 2314). This book writes its \emph{l} as \mw{ळ}, as in \textbf{\mw{काळो}} and \textbf{\mw{ओळखणो}}.
\end{cousinweb}

\begin{grammarlens}[title={What you've built}]
One more word for this chapter.
\end{grammarlens}

\subsection*{Your turn}
\begin{itemize}
\raggedright
  \item \emph{Say it:} \emph{oḷmo}
  \item \emph{Say it:} \emph{oḷmo}, once more
  \item \emph{Say it:} \emph{yo oḷmo sācho hai.} — This complaint is fair.
  \item \emph{Recall:} say the Marwadi for modesty, then the Marwadi for a sigh, then say \emph{oḷmo} again
\end{itemize}

\begin{tcolorbox}[breakable,colback=teal!4,colframe=teal!35!black,title={Before you move on}]
What does \mw{ओळमो} mean? (\textquotedblleft{}A reproach; a complaint\textquotedblright{}.) Say it once more.
\end{tcolorbox}
